10.5 Higher-Order Derivatives
83
Then, after computing the value of this operator at the vector $h_2 \in X$, we obtain an
element of $Y$:
$$f^{(2)} (x) (h_1, h_2) := (f^{(2)} (x)h_1)h_2 \in Y.$$
(10.65)
But
$$f^{(2)} (x)h = (f')'(x)h = D_h f'(x),$$
and therefore
$$f^{(2)} (x) (h_1, h_2) = (D_{h_1} f'(x))h_2.$$
(10.66)
If $A \in \mathcal{L}(X; Y)$ and $h \in X$, this pairing with $Ah$ can be regarded not only as a
mapping $h\mapsto Ah$ from $X$ into $Y$, but as a mapping $A \mapsto Ah$ from $\mathcal{L}(X; Y)$ into $Y$,
the latter mapping being linear, just like the former.
Comparing relations (10.62), (10.64), and (10.66), we can write
$$(D_{h_1} f'(x))h_2 = D_{h_1}(f'(x)h_2) = D_{h_1} D_{h_2} f (x).$$
Thus we finally obtain
$$f^{(2)} (x)(h_1, h_2) = D_{h_1} D_{h_2} f (x).$$
Similarly, one can show that the relation
$$f^{(n)}(x)(h_1,\ldots, h_n) := (\ldots (f^{(n)}(x)h_1)\ldots h_n) = D_{h_1} D_{h_2} \cdots D_{h_n} f (x) \quad (10.67)$$
holds for any $n \in \mathbb{N}$, the differentiation with respect to the vectors being carried
out sequentially, starting with differentiation with respect to $h_n$ and ending with
differentiation with respect to $h_1$.

10.5.3 Symmetry of the Higher-Order Differentials
In connection with formula (10.67), which is perfectly adequate for computation as
it now stands, the question naturally arises: To what extent does the result of the
computation depend on the order of differentiation?
Proposition If the form $f^{(n)} (x)$ is defined at the point $x$ for the mapping (10.58), it
is symmetric with respect to any pair of its arguments.
Proof The main element in the proof is to verify that the proposition holds in the
case $n = 2$.
Let $h_1$ and $h_2$ be two arbitrary fixed vectors in the space $T_Xx \sim X$. Since $U$ is
open in $X$, the following auxiliary function of $t$ is defined for all values of $t \in \mathbb{R}$
sufficiently close to zero:
$$F_t(h_1, h_2) = f (x + t(h_1+h_2)) - f (x + th_1) − f(x + th_2) + f (x).$$